\chapter{Boundary Radiation, Finite Edition}

The previous chapter generated a two-channel gauge from the boundary and left the two
channels as single actions. This chapter returns each channel to the family setting the
earlier towers used, and in doing so it earns a result the single-action picture could
not state: when the certificate family and the charge family are kept apart, and the
interior is closed against their mixing, the remaining charge has only one place to go.
It reads at the boundary, as a flux. The chapter is the finite edition of boundary
radiation, and every step of it is carried by an object the book has already built.

The chapter opens with a separation it makes by type rather than by theorem. The
certificate family and the electron family do not merely happen to be different; they
live in different sectors of a carrier that cannot mix them. The apparatus builds a
carrier with a geometric sector and a gauge sector and a pairing that is zero across the
sectors by construction --- not proven zero, made zero, the way a type forbids two things
from combining. There is no operation that adds a geometric value to a gauge value; the
two simply cannot be formed into one. The certificate family and the electron family are
orthogonal because they are, by type, unmixable.

The second section gives that orthogonality teeth with a piece of Clifford algebra in
shadow. Clifford multiplication remembers a pairing through its anticommutator: the
symmetric product of two elements is minus twice their pairing. So when two
representatives are orthogonal, the scalar part of their Clifford product vanishes, and a
nonzero interior mixed certificate --- a scalar coupling between the geometric and the
gauge sector --- would require changing the very structure that made the orthogonality
hold. Inside the unmodified apparatus there is nothing to write such a certificate with.
The interior is closed against mixing, conditionally but exactly: change the pairing, the
form, the representation, or the product, and a mixed certificate becomes writable --- and
that changed structure is what the book would call new physics.

The third section follows the only path the closed interior leaves open. If the charge
cannot appear as an interior mixed certificate, and yet the charge is real --- the
electron was produced --- then the charge must appear somewhere else, and the only
remaining place is the boundary. The apparatus reads the generated program's residue as a
boundary flux and proves it equals the terminal residual: a signed unit, the electron,
read at the boundary precisely because the interior refused it. This is boundary
radiation in the finite edition: a charge that the interior cannot hold, read as a flux
through the edge.

The fourth section pays a debt the earlier chapters left open. The produced tower was
built once, for a single showcase action. The apparatus now files it for every flat pair:
wherever the action is flat at a path, its residue vanishes identically, the tower's bound
is proven zero, and the completed tower exists with nothing assumed. The section is strict
about naming. A flat tower --- a produced tower over a flat pair --- is not a certificate of
admissibility, and the apparatus refuses to let one testify as the other. A flat tower
carries convergence; an admitted certificate carries admissibility; sliding from the first
to the second is a cheat the apparatus names and declines.

The fifth section discharges the book's oldest promise. The first chapter built the
machinery of a converging process long before anything could run it. Here it runs. The
gauge reading sits, at every stage, under the geometric bound the completed tower carries
--- and that bound, the geometric tolerance, tends to zero. So the gauge reading is
forced to be Cauchy: any two late readings differ by less than any tolerance, because both
sit under a geometric tolerance that has already vanished. The gauge converges because the
geometry's tolerance does, and it converges to the on-shell zero. The convergence is of
the strongest kind the book has: eventual exactness, a sequence that converges because it
has stopped moving.

By the end of the chapter the two families have been separated, the interior closed, the
charge read at the boundary, the tower filed for every flat pair, and the gauge reading
shown to converge as the geometric tolerance vanishes. The single invariant that became a
charge now radiates at the boundary, and the radiation converges because the geometry
spends its tolerance down to zero. Every claim is finite, and every claim is gated: this
is the finite shadow of boundary radiation, one signed unit at one boundary, and the
larger names keep their own references.

\section{The Split: Certificate Family versus Electron Family}

The chapter's first move is to separate two families so completely that they cannot be
combined, and to do the separating by type rather than by argument. The certificate family
--- the configurations the apparatus certifies through second order --- and the electron
family --- the charge that lives in the coupling --- are not two regions of one space that
happen not to overlap. They are two sectors of a carrier built so that nothing crosses
between them.

Recall where each family lives, because the separation reflects a real difference. The
certificate family lives where countable exhaustion suffices: below the door, where the
skeleton tower runs out of configurations and the residual becomes exactly zero. The
electron lives in the coupling that one-coordinate exhaustion cannot express --- the
strictly-between content that the weak form projects out and that only a paired variation
reveals. One family is reachable by exhausting a countable tower; the other is precisely
what such exhaustion leaves behind. They sit on opposite sides of a divide the countable
machinery cannot cross.

The apparatus encodes this as a two-sector carrier. A value is either geometric or gauge,
and there is no third option and no combination:
\[
  \text{value} = \text{geometric}(g)
  \quad\text{or}\quad
  \text{value} = \text{gauge}(q),
  \qquad\text{never both.}
\]
The pairing on this carrier respects the sectors. Two geometric values pair by the
geometric pairing; two gauge values pair by the gauge pairing; and a geometric value
paired with a gauge value is zero --- not because a computation returns zero, but because
the pairing is defined to be zero across sectors:
\[
  \langle \text{geometric}(g),\, \text{gauge}(q) \rangle = 0
  \qquad\text{by definition.}
\]
This is orthogonality supplied by type. The geometric and gauge sectors are orthogonal not
as a fact about a common space but as a rule about a carrier built so they cannot mix.

The sharpest expression of the separation is that the apparatus cannot even form the sum
of a geometric value and a gauge value. The carrier is a sum of sectors, not a direct sum
of spaces; there is no addition that takes a geometric value and a gauge value and returns
a combined value. One cannot write ``the geometric part plus the gauge part'' because the
operation does not exist. The impossibility of forming the sum is the point: the two
families do not coexist in a space where they could be added and their cross-term studied.
They live in separate sectors, and the separation is built into what a value is.

The representatives keep the split from being empty. The geometric sector carries a
nontrivial energy witness, which pairs with itself to a positive value; the gauge sector
carries the electron, which pairs with itself to the square of its charge. Each sector is
genuinely inhabited and genuinely nondegenerate, so the split is a separation of two real
families, not a partition of one populated sector from an empty one. The geometric witness
and the electron are both there, each carrying its own self-pairing, and the cross-pairing
between them is zero by type.

The section's metaphysical anchor is that some separations are matters of kind, not of
distance. The apparatus tanges the certificate family into a geometric sector and the
electron into a gauge sector, and it funges no value across the boundary between them
because the boundary is built into the carrier. The two families are orthogonal in the
strongest sense a finite apparatus can manage: not far apart in a common space, but
unmixable by construction, with no operation that could ever combine them. The chapter
will read the charge at the edge precisely because, by type, it has nowhere in the
interior to be.

\section{The Clifford Shadow and the Interior Obstruction}

The orthogonality of the previous section is, so far, a rule about a pairing. This section
gives it dynamical teeth through a fragment of Clifford algebra, read in shadow, and shows
that orthogonality closes the interior against a mixed certificate. The argument is short,
but it is the hinge on which the boundary reading turns.

Clifford multiplication is the algebra that remembers a quadratic form, and it remembers
it through its anticommutator. For two elements, the symmetric product --- the
anticommutator --- is minus twice their pairing:
\[
  c(u)\,c(v) + c(v)\,c(u) = -2\,\langle u, v \rangle.
\]
This relation is the whole of the Clifford structure the apparatus needs. It says that the
scalar part of a product of two elements is not free: it is determined by their pairing,
rescaled by minus two. The apparatus carries this relation as a shadow --- an
anticommutator that equals minus twice the pairing --- over any certificate split.

Apply it to the geometric and gauge representatives, which the previous section made
orthogonal. Their pairing is zero, so the scalar part of their Clifford product is zero:
\[
  \langle g, q \rangle = 0
  \quad\Longrightarrow\quad
  c(g)\,c(q) + c(q)\,c(g) = -2 \cdot 0 = 0.
\]
The scalar cross-term between the geometric sector and the gauge sector vanishes. This is
not an accident of the representatives; it follows from the orthogonality the type
supplies, through the Clifford relation, in one rewriting and one arithmetic step.

The consequence is an obstruction. Call a nonzero scalar cross-term between the two
sectors an interior mixed certificate --- a scalar coupling that would let the geometric
and the gauge sector interact inside the apparatus. Orthogonality forbids it: where the
pairing is zero, the scalar cross-term is zero, so no interior mixed certificate exists:
\[
  \langle g, q \rangle = 0
  \quad\Longrightarrow\quad
  \text{no interior mixed scalar certificate.}
\]
Inside the unmodified apparatus there is nothing to write such a certificate with. The
charge cannot couple to the geometry through a scalar in the interior, because the only
scalar available is the cross-term, and the cross-term is zero.

The obstruction is conditional, and the apparatus is careful to say so. It is not a
universal impossibility; it is an impossibility \emph{in the unmodified apparatus}. A
nonzero interior mixed certificate would require changing at least one of the structures
that made the orthogonality hold: the pairing, the form, the representation, or the
product. Change any of those and a mixed certificate becomes writable. That changed tuple
is exactly what the book would call new physics --- a different apparatus, with a different
interior, in which the geometry and the charge can couple inside. The current apparatus
forbids it; a modified one might allow it; and the conditional shape of the obstruction is
what makes it a statement about \emph{this} physics rather than a dogma about all possible
ones.

The section's metaphysical anchor is that orthogonality, read through the Clifford shadow,
closes a door. The apparatus tanges the two representatives into a Clifford product and
funges the pairing into its scalar part, and where the pairing is zero the scalar part is
zero, so the interior has no certificate to hold the coupling. The charge is real and the
interior is closed against it. There is exactly one place left for a real charge to be
read, and the next section reads it there: at the boundary.

\section{Boundary Flux Equals the Terminal Residual}

The previous section closed the interior against the charge. This section follows the only
path that leaves open. The charge is real --- the electron was produced, a signed unit
that no orthogonality can wish away --- and the interior cannot hold it. A real quantity
that the interior cannot hold does not vanish. It reads at the boundary.

State the situation as a syllogism, because its force is logical rather than computational.
The charge exists: the pair kick produces it. The interior is closed: orthogonality
forbids an interior mixed certificate. The charge must therefore be readable somewhere
other than the interior, and the apparatus has only two places, interior and boundary. The
interior is ruled out. The boundary remains:
\[
  \text{charge real} \ \wedge\ \text{interior closed}
  \;\Longrightarrow\;
  \text{charge reads at the boundary.}
\]
The boundary is not chosen for convenience. It is forced as the complement of a closed
interior.

Read there, the charge appears as a flux. The apparatus reads the generated program's
observable remainder --- the excitation channel's pair coupling --- as a flux through the
boundary, and it proves that this boundary flux equals the terminal residual:
\[
  \text{boundary flux} = \text{terminal residual} = \text{electron} = -1.
\]
The signed unit that the interior refused appears, undiminished, as a flux through the
edge. It is the same electron of the pair-production chapter, read now not in the interior
coupling but as radiation across the boundary, because the interior gave it nowhere to be.

This is the finite edition of boundary radiation, and the apparatus packages it as three
clauses, each carried by an object already filed. First, the geometric and gauge
representatives are orthogonal in the split pairing. Second, that orthogonality leaves no
interior mixed scalar certificate in the unmodified apparatus. Third, the observable
remainder reads as the boundary residual --- the electron, radiated. Orthogonality, closed
interior, flux equals residual: the three clauses of boundary radiation, finite edition,
each checked.

The apparatus is exact about the size of the claim, because boundary radiation is a phrase
with large associations. What is established here is one signed unit, read at one boundary,
in a finite apparatus. The orthogonality is supplied by the type of the carrier, not
proven about a continuum pairing; the obstruction is conditional on the apparatus being
unmodified; and the equality of flux and residual is verified at the instance, with the
general implication --- orthogonality and a closed interior forcing flux to equal residual
--- still a target rather than a discharged theorem. The larger phenomena that the words
``boundary radiation'' might summon keep their own references. The apparatus claims the
finite shadow: a charge the interior cannot hold, read as a flux at the edge.

The reading also explains, in retrospect, why the charge was always going to be a boundary
phenomenon. The charge lives in the coupling that the weak form projects out of the
interior; the interior, by the orthogonality, has no scalar to hold it; so the only place
the charge can be read is where the interior ends. A quantity defined by what the interior
discards is, of necessity, read at the interior's edge. Boundary radiation is not a
surprising place for the charge to appear. It is the only place left once the interior is
closed.

The section's metaphysical anchor is that what cannot be held inside is read at the edge.
The apparatus tanges the charge against a closed interior and funges it outward to the
boundary, where it appears as a flux equal to the terminal residual. The electron, refused
by the interior, radiates --- not because the apparatus chose to look at the boundary, but
because a real charge that the interior cannot hold has nowhere else to be read. Boundary
radiation, in the finite edition, is the charge appearing where the interior ends.

\section{The Flat Tower, Filed}

The chapter pauses here to pay a debt the earlier chapters left open, because the boundary
reading rests on a convergence that was only ever built for one case. The produced tower
--- the concrete, exhausting tower with its convergence proven rather than assumed --- was
constructed for a single showcase action. This section files it for every flat pair, and
guards a naming discipline that the rest of the book depends on.

The construction generalizes with one observation. Where an action is flat at a path ---
where the path is a critical point, the action stationary under every variation --- the
residual vanishes identically, because the residual is the derivative and a flat pair's
derivative is zero:
\[
  \text{action flat at path}
  \;\Longrightarrow\;
  \text{residual}(v) = 0 \ \text{ for every variation } v.
\]
With the residual identically zero, every stage of the tower carries a terminal residual
of zero, the squeeze's bound is proven to be zero rather than assumed to tend to zero, and
the completed tower exists with nothing assumed anywhere in it. The showcase construction
was not special; it worked because the showcase action was flat, and every flat pair
inherits the same construction.

So the apparatus files the produced tower for every flat pair. For any action flat at any
path, there is a completed tower whose terminal residual magnitude tends to zero at every
variation --- and tends to zero in the strongest way, by being exactly zero from the start.
The convergence the boundary reading leaned on is now available not for one action but for
the whole class of flat pairs, each carrying its own produced tower.

The section's real work, though, is a naming discipline, and the apparatus is strict about
it because the temptation to blur it is real. A flat tower and an admitted certificate are
different objects, and the apparatus refuses to let one testify as the other:
\[
  \text{flat tower} : \text{convergence over a flat pair, no admissibility claimed},
\]
\[
  \text{admitted certificate} : \text{the } C^2 \text{ certificate, whose first clause is admissibility}.
\]
A flat tower says only that an action's residual converges along the skeleton; it makes no
claim that the action admits any path. An admitted certificate says that the configuration
is admissible and certified through second order. The first is a statement about
convergence; the second is a statement about a law holding. Sliding from ``this tower
converges'' to ``this law is admitted'' is a cheat --- it smuggles admissibility out of a
convergence that never claimed it --- and the apparatus names the two channels distinctly so
that the slide cannot happen silently.

The discipline matters because the flat case is exactly where the slide is tempting. A flat
action is stationary everywhere; its residual is zero everywhere; it converges trivially.
One could mistake that trivial convergence for the action being a certified law. It is not.
Flatness is a property of the residual, not a certificate of admissibility, and a flat
tower carries the former without the latter. The apparatus files the convergence and
withholds the admissibility, keeping the two certificates in separate names so that the
convergence cannot be cashed as a law it never earned.

The section's metaphysical anchor is that a convergence is not a law. The apparatus tanges
each flat pair into a produced tower and funges its residual to zero, filing the
convergence for the whole class --- but it guards the boundary between convergence and
admissibility, refusing to let a tower that merely converges testify that a law is
admitted. The flat tower is filed, and the cheat is named and declined. What converges has
converged; what is admitted is a separate certificate, and the apparatus does not let the
one wear the other's clothes.

\section{Cauchy Convergence under a Vanishing Geometric Tolerance}

The book's first chapter built the machinery of a converging process --- a residue, an
origin, and a limit away from the origin --- many chapters before anything could run it. It
said as much at the time: the pieces were there, waiting. This section runs them, and the
process that converges is the gauge reading, driven to its limit by the geometry.

The two programs are coupled through exactly one quantity: the bound. The gauge reading,
at every stage, sits under the geometric completion's bound --- the residual the apparatus
reads on the gauge side is no larger than the tolerance the geometric refinement grants it:
\[
  \text{gauge residual}(n) \;\le\; \text{geometric bound}(n).
\]
And the geometric bound tends to zero; it was proven to, not assumed, when the tower was
produced. This geometric bound is the tolerance the geometry spends --- the apparatus calls
it the geometric tolerance, the finite shadow of the small parameter a continuum theory
would send to zero. The gauge reading is leashed to it.

A reading leashed to a tolerance that vanishes is forced to be Cauchy. Take any tolerance.
The geometric bound is eventually below it, so past some stage both of any two gauge
readings sit under the bound and hence under the tolerance, and so they differ from each
other by no more than the tolerance:
\[
  m, n \ge N
  \;\Longrightarrow\;
  \big|\,\text{gauge residual}(m) - \text{gauge residual}(n)\,\big| \le \varepsilon.
\]
For every tolerance there is such a stage. The gauge reading is Cauchy --- not by its own
internal effort, but because it is bounded by a geometric tolerance that has already gone
to zero. And it converges, to the on-shell zero: the limit is the stationary reading the
earlier chapters reached, approached here not by creeping but by arriving, since the
readings are eventually exactly zero. The gauge converges because the geometry's tolerance
does.

This is the discrete completion of the book's oldest promise, and it is honest about being
discrete. The tolerance is the finite, order-theoretic geometric bound, not a continuum
metric; the convergence is eventual exactness, the strongest Cauchy behavior a sequence can
have, which here has the flavor of a joke --- the sequence converges because it has stopped
moving. And Cauchy completeness, the existence of the limit inside a completed space,
remains behind the door the earlier chapters marked. What runs here is the finite shadow of
the converging process, with the geometry spending the tolerance and the gauge reading
following it down.

The section's metaphysical anchor is that one process can converge on another's account.
The apparatus tanges the gauge residual under the geometric bound and funges the geometry's
vanishing tolerance into the gauge reading's convergence, so that the charge's reading
settles because the geometry's tolerance settles. The residue the book named on its first
pages --- an origin and a limit --- reaches its limit here: the origin is the on-shell zero,
the limit is arrived at exactly, and the distance that shrinks is the geometry's to spend.
The gauge converges because the geometry's tolerance vanishes, and the oldest promise is
kept in the only currency the apparatus carries.

\section*{Bridge: Hygiene Before the Continuum Door}

The chapter separated two families by type, closed the interior against their mixing, read
the charge at the boundary, and let the geometry drive the gauge reading to its limit.
Every step was finite. But the separation it leaned on --- countable exhaustion reaching the
certificate family and stopping short of the coupling --- has the smell of a set-theoretic
divide, and the apparatus has been treating that divide informally. Before it approaches
the continuum in earnest, it owes a finite accounting of the separation it has been using.

The debt is specific. The certificate family was reached by exhausting a countable tower;
the electron was what such exhaustion could not reach. The apparatus has been calling these
opposite sides of a divide without saying, in finite terms, what the divide is and how its
two sides are kept apart. That is a set-theoretic question --- about counting, about
extension, about what a finite apparatus can decide and what it must leave open --- and the
apparatus has so far answered it by type rather than by an explicit finite hygiene.

The next chapter pays that debt with what it calls a finite forcing, a Cohen-style
construction read up to a tolerance. It counts the small structures the apparatus actually
distinguishes, files the finite spline conditions exactly, and separates what the apparatus
can extend and decide from what it must leave to the continuum door. The point is not to
prove a theorem of set theory but to keep the apparatus's own house in order: to say, in
finite terms, exactly how the two families are counted and kept apart before any continuum
claim is made.

So the next chapter files the set-theoretic hygiene. It counts to three in the apparatus's
own number system, states the finite spline conditions, sorts extension from compatibility
from decision, and closes out the gauge commutator residue. The question the apparatus
carries forward is the bookkeeping question its boundary reading raised: the two families
are kept apart, but by what finite counting, and how much of the keeping-apart is the
apparatus's to decide rather than the continuum's?

\section*{Coda: Oil, Water, and the Skin Between}

A drop of oil in a glass of water performs this chapter's whole argument without a word.
The two liquids do not mix; there is no interior in which they combine; and everything that
passes between them passes at the skin where they meet. The kitchen has been doing boundary
radiation, in its way, all along.

Begin with the refusal to mix, because it is a refusal of kind. Oil and water do not fail
to mix because no one has stirred hard enough. They fail to mix because of what they are ---
one polar, one not --- and no amount of stirring makes a homogeneous oil-water substance.
Stir vigorously and you get an emulsion, a suspension of droplets, but inside each droplet
the oil is still oil and around it the water is still water. There is no phase that is
genuinely both. The immiscibility is a matter of type, not of effort, exactly as the
certificate family and the electron family are unmixable by the carrier's type rather than
by any distance between them.

Because there is no interior mixing, there is no interior place for the two to interact.
Look inside the body of the oil and you find only oil; look inside the water and you find
only water. Neither carries a region where the two are coupled. If the oil and the water
are going to do anything to each other --- exchange a molecule, hold a tension, pass a
charge --- there is exactly one place it can happen: the interface, the thin skin between
them where surface tension lives. The interior is closed; the interaction is forced to the
boundary. That is the chapter's obstruction and its boundary reading, rendered in a meniscus:
no interior coupling, so all the flux runs through the skin.

And the skin settles as the system quiets. Stop stirring, and the agitation dies down;
the droplets coalesce, the emulsion breaks, and the two liquids separate into clean layers
with one flat interface between them. As the disturbance --- the tolerance the system is
willing to carry --- shrinks toward zero, the configuration converges to the simplest one:
oil above, water below, one skin between, at rest. The convergence is driven by the
disturbance vanishing, not by the liquids striving, exactly as the gauge reading converges
because the geometric tolerance vanishes rather than by any effort of its own.

One honest note keeps the metaphor exact. You can force oil and water to mix --- add a
surfactant, a soap, that bridges the polar and the nonpolar --- and then a stable mixed
phase exists where none did before. But that is a changed system, with a new chemistry
added; in the plain glass of oil and water, with nothing added, the two refuse the interior
and meet only at the skin. The mixed phase is available only to a modified system, exactly
as an interior mixed certificate is available only to a modified apparatus --- a different
pairing, a different product, the chapter's conditional ``new physics'' in the form of a
drop of soap.

That is the chapter in a glass. Immiscible by type, like the two families; no interior
mixed phase, like the closed interior; all exchange at the skin, like the boundary flux;
settling to a clean interface as the agitation vanishes, like the Cauchy convergence under a
vanishing tolerance; and a surfactant as the modified chemistry that would let the interior
mix, like the changed apparatus that would write an interior certificate. The electron
radiates at the boundary for the same reason oil meets water only at the skin: what cannot
combine inside can still meet at the edge.
